% !TeX program = xelatex
% ============================================================
% 数模论文 LaTeX 模板(国赛 CUMCM 格式)
% 编译方式: XeLaTeX(推荐, 中文支持最好); 或 LuaLaTeX
% 依赖宏包: ctex, amsmath, amssymb, booktabs, graphicx,
%           algorithm, algorithmic, multirow, float
% 用法: 复制本文件到你的 论文/ 目录, 按注释填写内容,
%       每个子问题复制 5.1 小节的结构.
% ============================================================
\documentclass[11pt]{ctexart}

% ---------- 页面设置: 国赛常用 A4 上下左右 2.5cm ----------
\usepackage[a4paper, top=2.5cm, bottom=2.5cm, left=2.5cm, right=2.5cm]{geometry}

\usepackage{amsmath, amssymb}   % 数学
\usepackage{booktabs}           % 三线表
\usepackage{graphicx}           % 插图
\usepackage{multirow}           % 表格合并
\usepackage{float}              % [H] 图表位置
\usepackage{caption}            % 图表标题
\captionsetup{font=small, labelfont=bf}

% ---------- 算法环境 ----------
\usepackage{algorithm}
\usepackage{algorithmic}
\renewcommand{\algorithmicrequire}{\textbf{输入:}}
\renewcommand{\algorithmicensure}{\textbf{输出:}}

% ---------- 代码环境(附录用) ----------
\usepackage{listings}
\lstset{basicstyle=\ttfamily\small, breaklines=true, frame=single}

\title{\textbf{论文题目: 研究对象与方法}}
\author{参赛队员1 \quad 参赛队员2 \quad 参赛队员3}
\date{}

\begin{document}
\maketitle

% ============================================================
% 摘要: 最后写! 每问一句"针对…问题, 首先…, 建立…模型,
% 采用…算法求解, 得到…结果; 通过…检验, 表明模型…"
% ============================================================
\begin{abstract}
【第一段: 背景 + 总体思路。】
针对问题一: 【……】。
针对问题二: 【……】。
针对问题三: 【……】。
本文模型在【指标】上达到【数值】, 可为【场景】提供参考。

\noindent\textbf{关键词:} 【词1; 词2; 词3】
\end{abstract}

\section{问题重述}
\subsection{问题背景}
【用自己的话压缩题目背景, 不照抄原文。】

\subsection{问题提出}
\begin{itemize}
  \item \textbf{问题一:}【整理为明确的数学任务: 研究对象、已知条件、任务】。
  \item \textbf{问题二:}【同上, 并说明与问题一的关系】。
  \item \textbf{问题三:}【同上】。
\end{itemize}

\section{问题分析}
【逐问: 题目要求 $\rightarrow$ 关键难点 $\rightarrow$ 数据/变量特征
$\rightarrow$ 总体建模思路 $\rightarrow$ 求解与检验方法。
写"为什么该模型适合本题", 不提前展开公式。】
\subsection{对问题一的分析}
【……】
\subsection{对问题二的分析}
【……注意承接问题一: 是否使用问题一的模型或输出。】
\subsection{对问题三的分析}
【……】

\section{模型假设}
\begin{enumerate}
  \item 【假设内容】【理由: 因为……所以可这样假设】;
  \item 【……】;
  \item 【……】.
\end{enumerate}
以上假设的合理性将在模型检验部分进一步验证。

\section{符号说明}
\begin{table}[H]
  \centering
  \begin{tabular}{clc}
    \toprule
    符号 & 含义 & 单位 \\
    \midrule
    $x_i$ & 第 $i$ 个…… & …… \\
    $\lambda$ & …… & …… \\
    \bottomrule
  \end{tabular}
  \caption{符号说明}
\end{table}

\section{模型建立与求解}

\subsection{问题一模型的建立与求解}

\subsubsection{问题一的分析与数据处理}
【题目要求 $\rightarrow$ 难点 $\rightarrow$ 数据特征 $\rightarrow$ 思路。
数据处理交代: 有什么问题、用什么方法、关键参数、为何这样选、
保留什么特征、效果展示。】

\subsubsection{【模型名】的建立}
【先文字说明思路, 再给公式, 再解释每一项含义。
公式前必须有引导文字("由……原理可得""考虑到……因素, 引入……")。】
\begin{equation}
  f(\mathbf{x}) = \sum_{i=1}^{n} w_i x_i \label{eq:model1}
\end{equation}
其中 $w_i$ 表示【……】, $x_i$ 表示【……】。式 \eqref{eq:model1} 用于计算【……】, 下一步将【……】。

% 优化问题的标准型
\begin{equation}
\begin{aligned}
  \max\ / \min \quad & f(\mathbf{x}) \\
  \text{s.t.} \quad & g_j(\mathbf{x}) \le 0, \quad j = 1, \dots, m, \\
  & \mathbf{x} \in D .
\end{aligned}
\label{eq:std}
\end{equation}

\subsubsection{模型求解}
【算法/伪代码/流程图。算法环境示例:】
\begin{algorithm}[H]
  \caption{【算法名】}
  \begin{algorithmic}[1]
    \REQUIRE 【输入: 数据、参数】
    \ENSURE 【输出: 最优解】
    \STATE 初始化种群/解集 $P_0$, 参数 $T_{\max}$;
    \FOR{$t = 1$ \textbf{to} $T_{\max}$}
      \STATE 【迭代更新规则: 如新解生成、适应度计算】;
      \STATE 【约束判断与修正】;
      \IF{满足终止条件}
        \STATE \textbf{break}
      \ENDIF
    \ENDFOR
    \RETURN 最优解
  \end{algorithmic}
\end{algorithm}

\subsubsection{结果与分析}
% 占位图: 换成你自己的图片; 文件不存在时自动跳过, 不影响编译
\IfFileExists{图1_xxx.png}{%
\begin{figure}[H]
  \centering
  \includegraphics[width=0.7\textwidth]{图1_xxx.png}
  \caption{图1 三种策略下的成本对比}
\end{figure}
}{}
【每张图/表必须在正文被引用并讨论("看图说话"); 对比表最优值加粗(\textbf{数值}); 结果是否符合预期与假设; 检验: 误差/残差/交叉验证/敏感性/多模型对比。】

\subsubsection{本问小结}
【"综上, 本文建立了……模型, 求得……, 模型在验证集上的指标为……, 说明模型……, 所得结果将在问题二作为……使用。"】

\subsection{问题二模型的建立与求解}
【同上结构。开头写承上启下过渡句: "问题二在问题一的基础上进一步考虑了……, 因此我们在模型一的基础上做了如下改进……"】

\subsection{问题三模型的建立与求解}
【同上结构。】

\section{模型评价与推广}
\subsection{模型优点}
\begin{itemize}
  \item 【……】
\end{itemize}
\subsection{模型缺点}
\begin{itemize}
  \item 【真实具体, 如"计算复杂度高, 数据量增大时求解时间超过……"】
\end{itemize}
\subsection{模型改进方向}
\begin{itemize}
  \item 【针对缺点给出可行改进思路】
\end{itemize}
\subsection{模型推广}
\begin{itemize}
  \item 【结合本题背景向外延伸, 不空谈】
\end{itemize}

\section{参考文献}
\begin{thebibliography}{9}
  \bibitem{ref1} 作者. 题名[J]. 期刊, 年, 卷(期): 页码.
  \bibitem{ref2} 作者. 题名[D]. 城市: 学校, 年.
\end{thebibliography}

\appendix
\section{详细推导}
【冗长推导放这里, 正文注明"详细推导见附录A"。】

\section{程序代码}
\begin{lstlisting}[language=Python]
# 核心代码
\end{lstlisting}

\end{document}
